\documentclass[11pt]{article}
\usepackage[a4paper,margin=30mm]{geometry}
\usepackage{amsmath,amssymb,amsthm,mathtools}
\usepackage[hidelinks]{hyperref}
\usepackage{microtype}

\newtheorem{theorem}{Theorem}[section]
\newtheorem{proposition}[theorem]{Proposition}
\newtheorem{lemma}[theorem]{Lemma}
\newtheorem{corollary}[theorem]{Corollary}
\theoremstyle{remark}
\newtheorem{remark}[theorem]{Remark}

\newcommand{\rhoA}{\rho}
\newcommand{\Z}{\mathbb Z}
\newcommand{\R}{\mathbb R}

\title{Periodic Flux Phases in Signed Circulants:\\Sharp Spectral Gaps and a Finite Threshold Transition}
\author{Working reconstruction draft --- author metadata to be confirmed}
\date{9 September 2026}

\begin{document}
\maketitle

\begin{abstract}
For $2\le s<N/2$, let
\[
 m(N,s)=\min_{\sigma}\rho(A_\sigma)
\]
be the minimum adjacency spectral radius over all edge signings of the
four-regular circulant $C_N(1,s)$.  A recent conjecture for $C_n(1,2)$
proposed a natural twisted signing as the global minimizer.  We first show
that this fails on an infinite family: a primitive period-eight flux phase
beats the twisted value whenever $8\mid n$ and $n\ge32$.

The counterexample is the first member of a general parity-dependent
mechanism.  For every jump $s\ge2$ we construct an explicit periodic signing
whose continuous squared Bloch spectral radius $\widehat R_s$ satisfies
$\widehat R_s<8$.  Writing $\widehat g_s=8-\widehat R_s$, we prove the
uniform estimates
\[
 \frac{1}{6s(s+2)}\le \widehat g_s
 \le 4\sin^2\!\frac{\pi}{s+2}
\]
and the sharp leading asymptotic law
\[
 s^2\widehat g_s\longrightarrow\pi^2.
\]
For even jumps the maximizing Bloch phase undergoes a boundary-layer phase
slip; its first rigorously audited correction occurs below the leading
$s^{-2}$ scale.

We also obtain two finite-graph classification results.  First,
$m(N,s)=2$ if and only if $N=2s+2$, while $m(N,s)\ge\sqrt5$ off this flat
line.  Second, on the resonance line $N=3s$,
\[
 m(3s,s)<\sqrt8
 \quad\Longleftrightarrow\quad
 s\ \text{is even or }s\in\{3,5\},
\]
and every odd $s\ge7$ obeys the uniform obstruction
$m(3s,s)^2\ge8+1/70$.  The exactly solvable case $s=2$ is used to explain
the half-cell chiral mechanism, the first sub-$\sqrt8$ primitive period, and
its rigidity.
\end{abstract}

\noindent\textbf{Keywords:} signed graph; circulant graph; spectral radius;
switching; Floquet decomposition; chiral symmetry; spectral gap

\medskip
\noindent\textbf{MSC 2020:} 05C50 (primary); 05C22, 15A18 (secondary)

\section{Introduction and main results}

Let $C_N(1,s)$ be the graph on $\Z/N\Z$ with edges at cyclic distances $1$
and $s$.  A signing $\sigma\in\{\pm1\}^{E}$ determines a real symmetric
signed adjacency matrix $A_\sigma$.  The fixed-support extremal problem is to
understand $m(N,s)$.

The introduction of the full manuscript should make the following logical
chain explicit:
\[
 \text{twisted-optimizer conjecture at }s=2
 \;\Longrightarrow\;
 \text{period-eight counterexample}
 \;\Longrightarrow\;
 \text{periodic flux mechanism for all }s
 \;\Longrightarrow\;
 \text{sharp }\pi^2/s^2\text{ gap}
 \;\Longrightarrow\;
 \text{finite threshold classification}.
\]

\begin{theorem}[Infinite counterexample family at jump two]
For every integer $L\ge4$, the graph $C_{8L}(1,2)$ has an explicit
period-eight signing $\sigma_L^*$ whose spectral radius is strictly smaller
than that of the natural twisted signing proposed as the global optimizer.
Consequently the twisted-optimizer conjecture fails for infinitely many even
orders.
\end{theorem}

\begin{theorem}[Periodic sub-threshold phases for every jump]
For every integer $s\ge2$ there is an explicit periodic signing of the
step-$(1,s)$ operator whose continuous squared Bloch spectral radius
$\widehat R_s$ is strictly smaller than $8$.  One may take a period-two
alternating-flux phase when $s$ is odd and a primitive period-$4s$
antipodal-defect phase when $s$ is even.
\end{theorem}

\begin{theorem}[Sharp periodic gap]
Let $\widehat g_s=8-\widehat R_s$ for the preceding explicit family.  Then
\[
 \frac{1}{6s(s+2)}\le \widehat g_s
 \le4\sin^2\frac{\pi}{s+2},
 \qquad
 s^2\widehat g_s\to\pi^2.
\]
For even $s=2r$, the global maximizing Bloch phase has the audited
boundary-layer asymptotics recorded in the quadratic-gap theorem package.
\end{theorem}

\begin{theorem}[Exact flat minimum]
For every admissible pair $(N,s)$,
\[
 m(N,s)=2\quad\Longleftrightarrow\quad N=2s+2.
\]
If $N\ne2s+2$, then $m(N,s)\ge\sqrt5$.
\end{theorem}

\begin{theorem}[Resonance-line threshold transition]
For every integer $s\ge2$,
\[
 m(3s,s)<\sqrt8
 \quad\Longleftrightarrow\quad
 s\text{ is even or }s\in\{3,5\}.
\]
Moreover, every odd $s\ge7$ satisfies
\[
 m(3s,s)^2\ge8+\frac1{70}.
\]
\end{theorem}

\section{Switching coordinates and periodic fibers}

This section is to be rebuilt from the reusable Hamilton-gauge and Bloch
material.  It should contain only notation and results needed later:

\begin{itemize}
  \item switching equivalence and cycle-flux coordinates;
  \item local flux versus global Hamilton holonomy;
  \item periodic cells and finite Bloch fibers;
  \item the general parity-dependent half-cell anticommutation criterion.
\end{itemize}

The period-eight-specific notation should not be introduced here unless it is
a specialization of the general $s$-jump setup.

\section{Periodic phases for arbitrary jump}

\subsection{Odd jumps}
Present the alternating-flux phase, diagonalize the Bloch symbol, and derive
the exact one-variable maximization problem.

\subsection{Even jumps}
Present the primitive period-$4s$ antipodal-defect phase.  Give the
Chebyshev/transfer reduction and prove that its squared Bloch edge is below
$8$ for every even $s$.

\subsection{A common spectral scale}
Explain why the parity-dependent constructions nevertheless share the same
leading gap scale.

\section{Sharp gap asymptotics and the even phase slip}

The proof order should be:

\begin{enumerate}
  \item coarse two-sided gap estimates;
  \item localization of the maximizing Bloch phase;
  \item the limit $s^2\widehat g_s\to\pi^2$;
  \item the even-jump local implicit equation;
  \item the rigorously audited second-order phase-slip correction.
\end{enumerate}

Exploratory critical-ratio calculations are excluded from theorem statements.

\section{Finite circulants: flat minima and resonance}

\subsection{The flat line}
Give the exact proof of $m(N,s)=2$ if and only if $N=2s+2$, followed by the
universal $\sqrt5$ lower bound off the flat line.

\subsection{The line $N=3s$}
Treat even $s$ by an explicit finite Fourier signing, $s=3,5$ by exact
positive-definiteness certificates, and odd $s\ge7$ by the global obstruction.
The finite scripts must appear only after the mathematical reduction proves
that the checked domain is complete.

\section{The exactly solvable base case $s=2$}

Retain the following period-eight results because they explain the general
construction:

\begin{itemize}
  \item half-cell flux/chiral equivalence;
  \item the exact $8\to4\to2$ fiber reduction and four squared bands;
  \item period eight is the smallest primitive period with edge below $8$;
  \item the first sub-eight period has a unique orbit under the natural
        dihedral/switching symmetries.
\end{itemize}

Long historical certificate tables and unrelated all-even computations should
move to supplementary material.

\section{Concluding problems}

The principal unsolved problem is
\[
 \boxed{\text{determine }m(N,s)\text{ exactly for every admissible }(N,s).}
\]
A second problem may ask for the minimal primitive sub-eight period and the
corresponding equality classes for general $s$.

\paragraph{Proof-source note for the reconstruction.}
The analytic and exact finite ingredients are presently distributed across
`research/paper_strengthening/`,
`research/generalization/circulant_1s/extension_20260905/`, and
`research/generalization/circulant_1s/quadratic_gap_20260906/`.
This skeleton is an article reconstruction, not a claim that those proofs have
already been merged and independently re-audited in this file.

\end{document}
